\documentclass[11pt,UTF8,fontset=fandol]{ctexart}
\usepackage{amsmath,amssymb,amsfonts}
\usepackage{graphicx,float,caption,subcaption}
\usepackage{booktabs,array,multirow,makecell}
\usepackage{geometry,enumitem,xcolor}
\usepackage[hidelinks]{hyperref}
\usepackage{url}
\geometry{margin=2.3cm}
\graphicspath{{/Users/yye/project/论文/assets/digirl/}}
\captionsetup{font=small,labelfont=bf}
\setlength{\parskip}{0.4em}

\newcommand{\ennote}[1]{\textit{\small(#1)}}

\title{\textbf{DigiRL：用自主强化学习训练\\“真实环境”里的设备控制智能体}}
\author{精读导读（中文讲解版）\\[2pt]
原论文：\textit{DigiRL: Training In-The-Wild Device-Control Agents with Autonomous Reinforcement Learning}\\[2pt]
Hao Bai, Yifei Zhou, Mert Cemri, Jiayi Pan, Alane Suhr, Sergey Levine, Aviral Kumar\\
UC Berkeley \quad UIUC \quad Google DeepMind}
\date{NeurIPS 2024 \quad|\quad 引用数 177（Semantic Scholar, 2026-06-07）\quad|\quad 代码：\url{https://github.com/DigiRL-agent/digirl}}

\begin{document}
\maketitle

\noindent\textbf{这篇在解决什么（一句话定位）：}
让一个不算大的视觉语言模型（VLM）学会“真刀真枪”地操作安卓手机界面（点击、滑动、打字），办法是让它\textbf{自己在真实安卓环境里反复尝试、靠成功/失败的反馈用强化学习（RL）变强}，而不是只靠人类示范去模仿。结果一个 1.3B 的小模型把任务成功率从 17.7\% 提到 67.2\%，反超了当时用 GPT-4V、Gemini 1.5 Pro 包装的大模型智能体，而且\textbf{全程不需要额外人工标注}。

\medskip
\hrule
\medskip

\section*{读前小白知识卡（先把这些概念搞懂）}
\addcontentsline{toc}{section}{读前小白知识卡}

\begin{itemize}[leftmargin=1.4em]
\item \textbf{智能体（agent）}：能感知环境、自己决定下一步动作、并对结果负责的程序。这里的“环境”就是一部安卓手机，“动作”就是点哪儿、滑哪儿、打什么字。
\item \textbf{VLM（视觉语言模型，vision-language model）}：既能看图、又能读懂/生成文字的大模型（如 GPT-4V、Gemini）。本文的智能体大脑就是一个 VLM：输入是手机截图＋任务指令，输出是一个动作。
\item \textbf{GUI（图形用户界面，graphical user interface）}：就是我们平时用的带图标、按钮、输入框的界面。“设备控制”指直接在 GUI 上按像素坐标点点划划，而不是调 API。
\item \textbf{In-the-wild（真实/野外环境）}：不是实验室里固定不变的模拟器，而是真实的 App 和网站——它们会改版、会弹广告、会让你登录、搜索结果顺序还会变。这种\textbf{随机性（stochasticity）}和\textbf{非平稳性（non-stationarity，环境随时间在变）}正是本文要对付的难点。
\item \textbf{模仿学习 / SFT（监督微调，supervised fine-tuning）}：拿人类操作的“标准答案轨迹”让模型照着学。问题：人类数据是\textbf{过时的（stale）}，而且模型一旦走到示范里没出现过的“没见过的状态”就不知道怎么办、爬不出来。
\item \textbf{强化学习（RL，reinforcement learning）}：让模型自己去试，做对了给奖励（reward），做错了不给，靠这个信号逐步学会更好的策略（policy）。本文的奖励很简单：整件任务\textbf{做成了给 1，没做成给 0}。
\item \textbf{在线 / 离线 RL（online / offline RL）}：\emph{在线}＝模型一边实际操作环境一边学（数据是它自己当下产生的）；\emph{离线}＝只拿一批事先收集好的旧数据来学，不再实际交互。本文是\textbf{先离线、再“离线转在线”（offline-to-online）}两段式。
\item \textbf{轨迹（trajectory / rollout）}：智能体从开始到结束完成（或没完成）一个任务，中间走过的“状态—动作”序列。一条轨迹就是一次完整尝试。
\item \textbf{信用分配（credit assignment）}：一条轨迹有很多步，最后只有一个总的成败信号。到底是哪一步走对了/走错了导致了成功/失败？把功劳或锅分到具体某一步上，就是信用分配问题。
\item \textbf{优势函数（advantage）}：衡量“在某个状态下做了某个动作，比平均水平好多少”。优势为正＝这个动作值得鼓励。本文核心就是\textbf{挑出优势大的动作让模型去学}（advantage-weighted）。
\item \textbf{价值函数（value function）}：估计“从某状态出发、按当前策略走下去，最终能拿到多少回报”。本文用了两个：\emph{指令级}（判断整条任务值不值得学）和\emph{步级}（判断某一步好不好）。
\item \textbf{AWR（优势加权回归，advantage-weighted regression）}：一种简单又稳的离线 RL 算法。直觉：把策略往“优势大的动作”那边拉，同时别离原来的策略太远。本文以它为骨架。
\end{itemize}

\medskip
\hrule
\medskip

\section{摘要全译}
视觉语言模型（VLM）的训练语料通常缺少足够的、以“决策”为中心的数据。这导致现成的 VLM 在决策类任务上表现欠佳——比如通过图形用户界面（GUI）在真实环境中控制设备。虽然“用静态示范数据训练”已显示出一定潜力，但我们指出：这类方法不足以控制真实的 GUI，因为它们无法应对真实世界中的随机性与非平稳性——而这些是静态观测数据里没有体现的。本文提出一种新颖的自主强化学习方法，称为 \textbf{DigiRL}，用于训练真实环境下的设备控制智能体；它通过两个阶段微调一个预训练 VLM：先用离线 RL 来初始化模型，再做“离线转在线”RL。为此，我们搭建了一个可扩展、可并行的安卓学习环境，并为其配备了一个基于 VLM 的评估器（evaluator）；同时为这一领域设计了一套简单却有效的 RL 方法。我们的方法运行\textbf{优势加权 RL}，其中的优势估计器经过增强以应对随机性，并配有一个\textbf{自动课程（automatic curriculum）}以榨取最大的学习信号。我们在 Android-in-the-Wild（AitW）数据集上验证了 DigiRL：用 RL 训练的 1.3B VLM 取得了 \textbf{49.5\% 的绝对提升}——成功率从 17.7\% 升到 67.2\%——超过了用静态人类示范数据做监督微调的版本。这一结果不仅显著超越此前最好的智能体（包括用 GPT-4V 的 AppAgent，成功率 8.3\%；以及用 AitW 数据训练的 17B CogAgent，38.5\%），也超过了此前最好的、基于“过滤式行为克隆”的自主 RL 方法（57.8\%），从而为真实环境下的设备控制数字智能体确立了新的最先进水平（SOTA）。

\section{逐节精读}

\subsection{引言：为什么“包装大模型”和“纯模仿”都不够（第 1 节）}

\textbf{愿景。} 作者想要的是一个真正自主的数字助手：能听懂自然语言、自己在电脑/手机上替你订机票、调 PPT、操作各种 App——而且要能在充满随机性和干扰的真实环境里稳定干活。

\textbf{核心痛点：会“推理”不等于会“干活”。} 哪怕是最强的闭源 VLM（GPT-4V、Gemini 1.5 Pro），在设备上执行任务时也经常给出错误动作。它们在抽象推理上很强，但这种能力\textbf{迁移不到“低层设备控制”上}。所以以往大多数做设备智能体的工作，都是给闭源 VLM 套上一层复杂的“外壳（wrapper）”：加提示词、加搜索、加工具调用。问题在于：\textbf{只要不更新模型权重，外壳能力再花哨，最终天花板也被底座模型卡死}。文中给的反例：让 GPT-4V “去 newegg.com 搜 alienware area 51 并选第一个”，它却跑去 Google 搜了 “alien area 51 ebay”，还自以为完成了任务（见图 2、图 15）。

\textbf{另一条路——纯模仿（imitation/SFT）也不行。} 网页和设备是动态的，照着\emph{过时}的示范数据学，环境一变就失效；而且这种智能体\textbf{无法从自己的错误里爬出来}——一旦走进示范里没见过的状态就卡死。

\textbf{本文主张：让模型在真实设备上从“自己的经验”里学。} 但这条交互式学习路线有两个硬性要求（desiderata）：
\begin{enumerate}[leftmargin=1.6em]
\item 必须用\textbf{在线交互数据}，因为静态示范代表不了部署时的真实情况（网站经常改版、有干扰）。
\item 必须能从\textbf{模型自己产生的多轮交互数据}里学，而这里面\emph{很大一部分是失败的}。所以要有机制\textbf{自动挑出对的动作、滤掉错的}。
\end{enumerate}

\textbf{贡献。} 提出自主 RL 方法 DigiRL（即“给数字智能体用的 RL”）。两段式：先离线 RL 初始化，再离线转在线 RL 用真实交互数据继续微调。为支撑在线 RL，作者搭了一个\textbf{可并行的安卓环境}（最多同时跑 64 个真实安卓模拟器），配一个\textbf{基于 VLM 的通用评估器}（对照人类判断的平均错误率仅 2.8\%）来自动给奖励。最终：1.3B 模型在 AitW 上比此前 SOTA（18B CogAgent，38.5\%）高出约 28.7 个百分点（到 67.2\%），比此前最好的自主学习方法（Filtered BC）高 9\% 以上，并大幅超过套在 GPT-4V/Gemini 上的外壳方案（17.7\%）。作者称这是\textbf{第一个成功把“离线转在线自主 RL”做到设备控制 SOTA 的工作}。

\subsection{相关工作（第 2 节）}
\begin{itemize}[leftmargin=1.4em]
\item \textbf{多模态数字智能体。} 与只处理文本/代码的语言智能体不同，控制设备是\textbf{像素级、坐标动作空间}的，而且设备+互联网生态随机不可预测。以往要么给闭源 VLM 套规则外壳（不微调＝天花板低），要么用示范数据做模仿学习（过时数据＋无法纠错）。最接近本文的是“在自采数据上做过滤式模仿学习”的工作。本文是\textbf{第一个把自主、离线转在线 RL 扩展到设备控制}的。即便对比那些在简化网页环境（MiniWob++）跑在线 RL 的工作，本文\textbf{样本效率高约 1000 倍}（约 $10^3$ 条轨迹 vs.\ 约 $10^6$ 条），且面向真实 GUI。
\item \textbf{设备控制的环境。} 已有的模拟环境大多只为“评测”设计，任务范围窄、完全确定且静态，学不到设备控制需要的丰富技能；另一些虽然任务多但过度简化、迁不到真实世界。本文环境用\textbf{自主评估 + Gemini 1.5 Pro}，支持在\textbf{并行的真实安卓设备}上做多样、开放式任务，规模远超以往（以往多是单线程模拟器，撑不起大规模在线 RL）。
\item \textbf{给 LLM/VLM 用的 RL。} 多数已有工作只针对\textbf{单轮}任务（如偏好优化、推理）。而从专家示范优化单轮，对多步问题往往次优，尤其在高随机/非平稳下。已有的基于价值的多轮 RL 算法通常要维护一堆网络（Q 网络、价值网络、策略网络及其延迟目标），收敛慢、对超参敏感。本文则聚焦\textbf{找出能落地的少数关键设计}，做一个简单有效、可被实践者直接用的 RL 算法。
\end{itemize}

\subsection{问题设定与预备知识（第 3 节）}

\textbf{研究范围。} 基于像素与虚拟设备交互，聚焦\textbf{控制安卓设备}——这比以往只做网页导航的环境更难也更通用（浏览器只是设备里的一个 App）。

\textbf{一个回合（episode）怎么走：}模拟器从主屏开始 $\rightarrow$ 从预设的语言指令集里选一个任务 $\rightarrow$ 智能体操作模拟器去完成。每一步，智能体看到的\textbf{观测就是当前屏幕截图}；动作空间沿用前人工作，包括：按归一化 $(x,y)$ 坐标（0$\sim$1 相对屏幕尺寸）\textbf{点击/滑动}、\textbf{输入任意长度文本}、以及按 HOME/BACK/ENTER 等\textbf{特殊键}（见图 3）。训练与测试指令来自 AitW 的 General（找信息、用基础 App，如“In-n-Out 的菜单上有什么？”）和 Web Shopping（网购，如“去 newegg.com 搜 razer kraken 并选第一个”）两个子集。

\textbf{随机性带来的挑战。} 真实设备控制有模拟环境里没有的随机性：(1) 网站/应用频繁更新（非平稳，导致在线观测和旧的离线数据对不上）；(2) 各种不可预测的干扰（弹窗广告、登录请求、搜索结果顺序随机）；(3) 技术故障（网页没加载完、临时访问受限）。

\textbf{“冻结策略会退化”的实证（图 4）：}作者先用 6 月 1$\sim$3 日的数据训出一个好模型，然后在 6 月 7$\sim$11 日分别测“冻结不更新的策略”和“持续在线更新的策略”。结果\textbf{冻结策略的成功率随时间逐渐下滑}（因为网站变了），而\textbf{持续在线更新能顶住这种退化}——这正是“必须在线学”的直接证据。

\textbf{可靠、可扩展的在线 RL 设置。} 自主 RL 是“边收数据边训练”，要在随机环境里学好，就得有\textbf{实时数据收集流水线}收集足够经验来做梯度更新。单线程安卓模拟器因延迟做不到，所以作者\textbf{并行化}模拟器（带错误处理）。此外环境必须能给\textbf{奖励信号}：判断当前屏幕是否表示任务已完成。作者扩展了 Pan et al.\ \ennote{[26]} 的\textbf{端到端自主评估器}——它\emph{不需要}访问模拟器内部状态、也\emph{不需要}给每个任务手写规则（这点和那些“给每个任务人工写校验函数”的工作形成对比）。评估器骨干用 \textbf{Gemini 1.5 Pro}，用少量示例轨迹＋人工标注的成败标签来引导它判断新任务。一个评估器即可覆盖所有 AitW 任务，且与人类标注高度一致（\textbf{平均错误率 2.8\%}，见图 8）。

\subsection{DigiRL 方法：为设备控制量身改造的自主 RL（第 4 节）}

把设备控制建成一个\textbf{有限步长的马尔可夫决策过程（MDP，Markov decision process）}：$M=\{S,A,T,\mu_0,R,H\}$，用策略梯度求解。开局从初始分布 $\mu_0$ 采样初始状态 $s_0$ 和自然语言指令 $c$；\textbf{只有在最后}——如果评估器判定任务成功——才给奖励 1，否则给 0；当任务完成或步数超过上限 $H$ 时回合结束。\textbf{状态用最近两张截图表示}。

几个 RL 标准定义（逐符号）：
\begin{itemize}[leftmargin=1.4em]
\item \textbf{Q 函数} $Q^\pi(s_h,a_h,c)=\mathbb{E}_\pi\!\left[\sum_{t=h}^{H} r(s_t,a_t,c)\right]$：在第 $h$ 步状态下做动作 $a_h$、之后一直按策略 $\pi$ 走，预期能拿到的\textbf{长期总回报}。
\item \textbf{价值函数} $V^\pi(s_h,c)$：把 $Q^\pi(s_h,a_h,c)$ 对“从 $\pi$ 采样的动作 $a_h$”求平均，即\textbf{这个状态本身的平均好坏}。
\item \textbf{优势} $A^\pi(s_h,a_h,c)=Q^\pi(s_h,a_h,c)-V^\pi(s_h,c)$：某个具体动作\textbf{比平均水平好多少}。
\end{itemize}

\subsubsection*{4.1 骨架：用 AWR 做离策略 RL}
出发点是 \textbf{AWR（优势加权回归）}：把策略往“奖励诱导的指数化优势”那边回归，作为“优化策略梯度但又别离上一版策略太远”的近似：
\[
\arg\max_\pi\ \mathbb{E}_\nu\big[\log \pi(a\mid s,c)\cdot \exp\big(A(s,a,c)/\beta\big)\big]. \tag{4.1}
\]
其中 $\beta>0$ 是温度超参，$\nu$ 是过去经验的分布，$A(s,a,c)$ 是动作的优势。\textbf{直觉：}优势越大的样本，权重 $\exp(A/\beta)$ 越大，模型就越被往那个动作上拉。

为了\textbf{省掉调 $\beta$ 的麻烦}，作者改成对优势做“硬过滤（hard filtering）”——不是按 $\exp(A)$ 连续加权，而是\textbf{直接挑出优势达标的样本}，然后对它们做普通的最大似然（模仿）训练：
\[
L(\pi)=-\mathbb{E}_{\mathrm{filter}(\nu)}\big[\log \pi(a\mid s,c)\big]. \tag{4.2}
\]
\textbf{说人话：}先用“优势”这把筛子把好动作挑出来，再让模型像背标准答案一样去学这些好动作。问题是：优势通常靠在环境里跑蒙特卡洛（MC）采样来估，而设备生态太随机，\textbf{MC 估出来的优势方差很大}（同样的动作有时成有时败）。

\subsubsection*{4.2 用“双稳健”估计器得到可靠的优势（doubly-robust）}
为了在强随机下也能可靠分辨“好动作”，作者构造了一个\textbf{逐步优势估计器}（受双稳健估计启发）：
\[
A^{\mathrm{step}}(s_h,a_h,c):=\lambda^{H-h} r(s_H,a_H,c)+\big(1-\lambda^{H-h} r(s_H,a_H,c)\big)\big(V^{\mathrm{step}}(s_{h+1},c)+r(s_h,a_h,c)-V^{\mathrm{step}}(s_h,c)\big). \tag{4.3}
\]
逐符号拆解：
\begin{itemize}[leftmargin=1.4em]
\item $\lambda$ 是一个加权超参；$\lambda^{H-h}$ 表示\textbf{离结尾越近、越信任“最终真实结果”这一项}。
\item 第一项 $\lambda^{H-h} r(s_H,a_H,c)$ 是\textbf{高方差的蒙特卡洛项}（直接看这条轨迹最后成没成，受随机性影响大）。
\item 后半括号里的 $V^{\mathrm{step}}(s_{h+1},c)+r(s_h,a_h,c)-V^{\mathrm{step}}(s_h,c)$ 是\textbf{高偏差的自举项}（用学到的步级价值函数估，价值函数拟合不完美会有偏）。
\item 因为本问题\textbf{中间没有奖励}，这其实是 GAE（广义优势估计）的简化版，只用“下一步优势”和“最终步优势”两端来插值。
\end{itemize}
\textbf{核心直觉：}把“高方差但无偏”和“高偏差但低方差”两种估计加权融合，找到\textbf{方差与偏差的甜蜜点}。实现“步级硬过滤”时很简单：当 $A^{\mathrm{step}}(s_h,a_h,c)>1/H$ 时，就认为这一步\textbf{在朝目标前进}，保留下来。

\subsubsection*{4.3 用“指令级价值函数”做自动课程（automatic curriculum）}
只有 4.1 的 AWR 更新＋4.2 的稳健优势，在普通 RL 任务上也许够了，但作者发现对设备控制还不够。原因：任务难度差异极大，\textbf{在“模型已经很熟练的任务”上反复收数据，会拖累样本效率}。最有效的学习信号来自\textbf{让智能体去经历那些最有信息量的任务}。

于是设计\textbf{指令级价值函数} $V^{\mathrm{instruct}}(c)$，判断某条轨迹值不值得学：
\[
A^{\mathrm{instruct}}(s_h,a_h,c):=\sum_{t=h}^{H} r(s_t,a_t,c)-V^{\mathrm{instruct}}(c)=r(s_H,a_H,c)-V^{\mathrm{instruct}}(c). \tag{4.4}
\]
（等号成立是因为\textbf{只有结尾才有奖励}，所以那个求和就等于最终奖励 $r(s_H,a_H,c)$。）\textbf{直觉：}如果一条轨迹的 $A^{\mathrm{instruct}}$ 很大，说明 $V^{\mathrm{instruct}}(c)$ 很小——即这是个\textbf{本来很难、但智能体居然做成了}的宝贵经验，应该优先学（类似“优先经验回放/关卡回放”的思想）。训练时先筛出 $A^{\mathrm{instruct}}$ \textbf{最高的前 $p$ 比例}数据，再丢给 AWR（配 4.3 的稳健优势）。

\textbf{实现细节：}受“现代深度网络（如 Transformer）用\textbf{交叉熵损失}训价值函数比用均方误差更好”的发现启发，两个价值函数都用\textbf{基于轨迹奖励 MC 估计的交叉熵目标}来训：
\[
L(V^{\mathrm{traj}})=-\mathbb{E}_\nu\big[r(s_H,a_H,c)\log V^{\mathrm{traj}}(c)+(1-r(s_H,a_H,c))\log(1-V^{\mathrm{traj}}(c))\big]. \tag{4.5}
\]
\[
L(V^{\mathrm{step}})=-\mathbb{E}_\nu\big[r(s_H,a_H,c)\log V^{\mathrm{step}}(s_h,a_h,c)+(1-r(s_H,a_H,c))\log(1-V^{\mathrm{step}}(s_h,a_h,c))\big]. \tag{4.6}
\]
\textbf{说人话：}既然奖励只有 0/1，就把价值预测当成“这条任务/这一步最终会成功的概率”，用\textbf{二分类的交叉熵}来学。

\textbf{最终算法（图 5）：}先用收集到的轨迹按 (4.5)/(4.6) 训\textbf{指令级}和\textbf{步级}两个价值函数；训演员（actor，即策略）时，先用 (4.4) 在指令级筛出“值得学的任务”、再用 (4.3) 在步级筛出“走对的那几步”，最后只在\textbf{筛过的数据}上用 (4.2) 的 MLE（最大似然）损失更新策略。

\subsection{方法图解：DigiRL 总览}

图 1 是理解整套方法的关键。它把 DigiRL 拆成三块，请配合图看下面这段“数据怎么流动”的讲解。

\begin{figure}[H]\centering
  \includegraphics[width=0.95\linewidth]{cand_p02_k00.pdf}
  \caption{DigiRL 总览（原文图 1）。DigiRL 建立在一个已用海量网络数据预训练过的 VLM 之上（具备常识、推理、视觉定位等基础能力）。首先用离线 RL、拿“过时的任务专属数据”微调 VLM，先把它的“目标导向行为”引出来；随后让智能体与真实 GUI 交互，通过在线 RL 与自主性能评估\textbf{持续}提升表现。}
  \label{fig:overview}
\end{figure}

\textbf{图里在说什么（顺着箭头走）：}
\begin{itemize}[leftmargin=1.4em]
\item \textbf{Pretraining（预训练，左上虚线框）：}起点是一个在互联网级图文数据上预训练好的 VLM，得到“Pretrained Model”。它有常识和看图能力，但还\textbf{不会“干活”}。
\item \textbf{Step I: Offline RL（离线 RL，左下框）：}拿一批\emph{已有的、可能过时的}轨迹，用离线 RL 微调，把“Pretrained Model”变成“Offline Model”。图里的对勾/叉号表示用成败信号筛选——\textbf{这一步先让模型有个像样的起点}，不用真实交互，省时间。
\item \textbf{Step II: Online RL（在线 RL，右侧大框，整套闭环）：}这是核心循环——
  \begin{enumerate}[leftmargin=1.6em]
  \item 从任务集采样任务（图中例子“show me a cat video”）；
  \item \textbf{Online Model 并行执行任务}，产生一批轨迹（多个安卓模拟器同时跑）；
  \item \textbf{AutoEval（自动评估器）}给每条轨迹判成败、打奖励（红叉＝失败、绿勾＝成功）；
  \item 把这些\textbf{自动标注好的轨迹}回灌，用在线 RL 更新模型（蓝色回环箭头），然后再去执行——\textbf{如此往复，越用越强}。
  \end{enumerate}
\end{itemize}
\textbf{一句话抓住精髓：}整条流水线\textbf{没有人工标注}介入——奖励由 AutoEval 自动给，数据由模型自己产生，模型靠这个闭环不断从“自己的成功与失败”里学习，从而能跟上真实环境的变化。

\section{实验解读（第 5 节）}

\textbf{用什么数据/基准：}AitW（Android-in-the-Wild）的两个子集——\textbf{General}（找信息/基础 App，最多 10 步）和 \textbf{Web Shopping}（网购，最多 20 步）。General 用全部 545 个训练任务，测试用测试集前 96 个；Web Shopping 经过后处理（剔除需要验证码/登录的网站与操作），训练 438 个、测试 96 个。

\textbf{对比哪些基线：}
\begin{itemize}[leftmargin=1.4em]
\item \textbf{提示/检索类（套在闭源 VLM 上）}：Set-of-Marks（给可点元素叠数字编号，让 VLM 直接输出编号而非算坐标）＋ GPT-4V / Gemini 1.5 Pro；以及 AppAgent（先让 VLM 探索、把经验加进测试提示）。
\item \textbf{监督训练类}：CogAgent（18B，用 AitW 训）、AutoUI（约 1.5B：200M 语言模型＋1.1B 视觉编码器）。
\item \textbf{自主学习类}：Filtered BC（过滤式行为克隆，此前最强的自采数据方法）。
\end{itemize}
\textbf{DigiRL 自己的底座：}基于 AutoUI-Base，\textbf{冻结图像编码器}；价值函数用同一个冻结编码器的视觉特征，拼上 RoBERTa 抽的指令特征，再过一个两层 MLP 预测。为公平起见，纯离线训练用的数据量 $\approx$ 离线转在线训练的总数据量（两子集离线约 1500 条轨迹；离线转在线＝约 500 条离线起步 + 约 1000 条在线）。

\subsection*{主结果（表 1）：小模型大反超}

\begin{table}[H]\centering
\caption{不同智能体在各设置下的成功率（\%）对比（原文表 1，节选）。离线实验跑 3 次报均值±标准差，在线实验跑 2 次；用自主评估器在训练/测试集前 96 条指令上评测。}
\begin{tabular}{llcccc}
\toprule
\multirow{2}{*}{类别} & \multirow{2}{*}{方法} & \multicolumn{2}{c}{AitW General} & \multicolumn{2}{c}{AitW Web Shopping} \\
\cmidrule(lr){3-4}\cmidrule(lr){5-6}
 & & 训练 & 测试 & 训练 & 测试 \\
\midrule
\multirow{4}{*}{提示} & Set-of-Marks · GPT-4V & 5.2 & 13.5 & 3.1 & 8.3 \\
 & Set-of-Marks · Gemini 1.5 Pro & 32.3 & 16.7 & 6.3 & 11.5 \\
 & AppAgent · GPT-4V & 13.5 & 17.7 & 12.5 & 8.3 \\
 & AppAgent · Gemini 1.5 Pro & 14.6 & 16.7 & 5.2 & 8.3 \\
\midrule
\multirow{2}{*}{监督训练} & CogAgent (18B) & 25.0 & 25.0 & 31.3 & 38.5 \\
 & AutoUI ($\approx$1.5B) & 12.5 & 14.6 & 14.6 & 17.7 \\
\midrule
\multirow{2}{*}{学习·离线} & Filtered BC & $51.7\pm5.4$ & $50.7\pm1.8$ & $44.7\pm1.6$ & $45.8\pm0.9$ \\
 & \textbf{本文（离线）} & $46.9\pm5.6$ & $62.8\pm1.0$ & $39.3\pm6.0$ & $45.8\pm6.6$ \\
\midrule
\multirow{2}{*}{学习·离线转在线} & Filtered BC & $53.5\pm0.8$ & $61.5\pm1.1$ & $53.6\pm4.7$ & $57.8\pm2.6$ \\
 & \textbf{DigiRL（离线转在线）} & $\mathbf{63.5\pm0.0}$ & $\mathbf{71.9\pm1.1}$ & $\mathbf{68.2\pm6.8}$ & $\mathbf{67.2\pm1.5}$ \\
\bottomrule
\end{tabular}
\end{table}

\textbf{这些数字说明了什么：}
\begin{itemize}[leftmargin=1.4em]
\item \textbf{光靠提示/检索不够。} GPT-4V、Gemini 即便加了 Set-of-Marks，成功率也都\textbf{不到 20\%}；AppAgent 这类推理时适配带来的提升\textbf{不超过 5\%}。说明\textbf{不微调，提升空间很有限}。失败主要来自“幻觉式推理”——跑到一个相关但错误的页面还自以为对了。
\item \textbf{监督训练有用但触顶。} 专门小模型 AutoUI（1.5B）能追平甚至超过通用大模型，但两个子集成功率仍\textbf{低于 20\%}；把模型堆到 18B（CogAgent）也还\textbf{低于 40\%}——\textbf{这不是靠堆规模/改架构能根治的}。这些智能体最典型的毛病是\textbf{无法纠正自己的错误}（比如误点了语音输入按钮，到步数上限都退不出来）。
\item \textbf{在线交互是关键。} 即便离线和“离线转在线”用的样本总量一样，\textbf{离线转在线显著更强}（General 子集 75\% vs.\ 62.8\%）——证明\textbf{让模型实际去环境里交互、从未经筛选的次优数据里学}非常有效。
\item \textbf{DigiRL 全面超过 Filtered BC。} 在 General 从 61.5\%$\to$71.9\%、在 Web Shopping 从 57.8\%$\to$67.2\%。\textbf{综合看：1.3B 的 DigiRL 把成功率从 17.7\% 提到 67.2\%（绝对 +49.5），反超 18B 大模型和所有外壳方案，且全程零人工标注。}
\end{itemize}

\subsection*{消融与分析（第 5.2 节）}

\textbf{失败模式分析（图 7）。} 把失败分三类：(1) \textbf{无法从错误中恢复}（走错后爬不出来，如停在错的搜索页）；(2) \textbf{中途卡住}（在正确路上分心，比如没点对链接、打完字忘了搜）；(3) \textbf{到了错误目标}（跑到错页面却以为完成了）。三类都因 RL 训练而减少，但\textbf{下降最稳、最显著的是“无法从错误中恢复”}——因为预训练模型只会“顺着生成看似合理的下一步”，遇到没见过的状态就抓瞎；而 DigiRL 在\textbf{自采轨迹}上训练，学会了从自己的错误里恢复。

\textbf{各组件消融（图 9 左，在 Web Shopping 上）。} 四个设计都必要：
\begin{itemize}[leftmargin=1.4em]
\item 价值函数用\textbf{交叉熵}（而非回归/均方）：约 \textbf{+12\%}（Ours vs.\ Ours w/ Regression）。
\item 用\textbf{步级优势}：效率约 \textbf{+12\%}（vs.\ 去掉步级优势）。
\item \textbf{自动课程}：学习速度约 \textbf{+25\%}（去掉步级优势的版本 vs.\ Filtered BC 对照）。
\item 完整版优于\textbf{原版 AWR}（没有双稳健优势、没有课程）。此外\textbf{“硬过滤”不掉点}——和精调温度 $\tau$ 后的指数过滤几乎打平，但实现更简单。
\end{itemize}

\textbf{评估器的可靠性（图 8）。} 自动评估器报的成功率与人类评估\textbf{几乎一致，差异 $<3\%$}；Web Shopping 上更准（语言表达更规整、歧义少），General 偏口语化所以略难。

\textbf{并行模拟器的加速（图 9 右）。} 随机器数量\textbf{近线性加速}：分布式模拟器在 128 CPU 上可稳定跑 64 个并行模拟器；而“所有并行都挤在同一台服务器”的朴素基线差很多（64 CPU 下 0.74 vs.\ 1.74 轨迹/分钟）。

\subsection*{附录里的额外发现}
\begin{itemize}[leftmargin=1.4em]
\item \textbf{步长上限（附录 B.1）。} 给更“富余”的步长（$H{=}20$ vs.\ $10$）能让两种方法都学得更快（一条轨迹里试错机会更多）；两种设置下 DigiRL 都比 Filtered BC \textbf{快约 100 条轨迹}。
\item \textbf{轨迹更短（附录 B.2、表 4）。} DigiRL 的平均轨迹长度\textbf{一直比 Filtered BC 短}——说明它解题更高效、更“直奔目标”。
\end{itemize}

\subsection{讨论与局限（第 6 节）}
作者用 DigiRL 在 AitW 多个安卓控制任务上确立了新 SOTA：先搭可扩展、可并行、带稳健 VLM 评估器的安卓环境支撑快速在线收数据；再做“离线 RL 预训练 + 自主 RL 微调”，在真实互联网/设备生态的随机性中靠交互学习。相对此前 SOTA 实现了\textbf{约 280\% 的相对提升}（成功率从 17.7\% 到 68.2\%）。\textbf{局限：}受算力所限，目前只在 AitW 的任务上训练，而非设备上所有可能任务（尽管并行模拟器和自动评估器都能扩展到更复杂任务）；算法刻意追求\textbf{实现简单}，希望成为后续研究（算法层面、以及扩展任务空间）的\textbf{基线}。

\section{核心 takeaways（给 AI 应用开发者的启示）}
\begin{enumerate}[leftmargin=1.6em]
\item \textbf{“包装大模型”有天花板，闭环学习才有上限。} 给 GPT-4V/Gemini 套提示词/检索/工具，短期见效，但不更新权重，能力被底座卡死。要让智能体在垂直场景里真正变强，往往得让它\textbf{从自己的交互反馈里学}。
\item \textbf{真实环境会变，所以模型也得一直学。} 图 4 的“冻结策略退化”是个强警示：网站/App 一改版，离线训练好的智能体就开始失效。\textbf{生产级 agent 需要持续在线适配（continual/online adaptation），而不是一次训练终身使用。}
\item \textbf{用一个 VLM 当“自动评委”是性价比极高的工程套路。} 不写每任务的规则、不靠人标，用强 VLM（这里是 Gemini 1.5 Pro）＋少量示例就能给出与人类一致（错误率 2.8\%）的成败判断——\textbf{这把“RL 没有奖励信号”这个最大障碍解决了}，且可扩展到很多任务。做 LLM-as-judge / 自动评测时可借鉴。
\item \textbf{小模型 + 好方法 > 大模型 + 纯提示。} 1.3B 反超 18B 与各种外壳，关键在\textbf{两段式训练（先离线起步、再在线精炼）}和\textbf{把成败信号正确分配到具体步骤}。算力有限时，方法和数据闭环比单纯堆参数更划算。
\item \textbf{把工程难点也当成研究的一部分。} 能做在线 RL 的前提是\textbf{并行、实时、带容错的环境}（这里是 64 个并行安卓模拟器，近线性加速）。\textbf{数据收集速度往往才是真正的瓶颈}——基础设施做对了，算法才跑得起来。
\item \textbf{保持简单、可复现。} 作者特意不堆一大堆 Q/V/策略网络，而是用 AWR 这种“筛出好动作再模仿”的简单骨架，加两三个关键改造（双稳健优势、交叉熵价值、自动课程）。\textbf{对工程落地而言，简单稳健、超参不敏感，比花哨更重要。}
\end{enumerate}

\section*{参考文献（节选最关键几条，英文原样）}
\addcontentsline{toc}{section}{参考文献}
\begin{enumerate}[leftmargin=2em]
\item C. Rawles, A. Li, D. Rodriguez, O. Riva, T. Lillicrap. \textit{Android in the Wild: A Large-scale Dataset for Android Device Control}. arXiv:2307.10088, 2023.
\item J. Pan, Y. Zhang, N. Tomlin, Y. Zhou, S. Levine, A. Suhr. \textit{Autonomous Evaluation and Refinement of Digital Agents}. arXiv:2404.06474, 2024.
\item X. B. Peng, A. Kumar, G. Zhang, S. Levine. \textit{Advantage-Weighted Regression: Simple and Scalable Off-Policy Reinforcement Learning}. arXiv:1910.00177, 2019.
\item Z. Zhang, A. Zhang. \textit{You Only Look at Screens: Multimodal Chain-of-Action Agents (AutoUI)}. 2023.
\item W. Hong et al. \textit{CogAgent: A Visual Language Model for GUI Agents}. 2023.
\item J. Schulman, P. Moritz, S. Levine, M. Jordan, P. Abbeel. \textit{High-Dimensional Continuous Control Using Generalized Advantage Estimation (GAE)}. 2018.
\item J. Farebrother et al. \textit{Stop Regressing: Training Value Functions via Classification for Scalable Deep RL}. 2024.
\end{enumerate}

\end{document}
